% ============================================================================
% CHAPTER 10: EVALUATION, METRICS, AND DECODING
% Where theory meets reality: measuring what matters in NLP.
% ============================================================================

\chapter{Evaluation, Metrics, and Decoding}
\label{chap:evaluation}

\begin{tldr}
Evaluation is where theory meets reality.  This chapter covers the complete landscape of
NLP evaluation: language modeling metrics (perplexity and its connection to cross-entropy),
generation metrics (BLEU, ROUGE, METEOR, BERTScore, COMET), classification metrics
(precision, recall, F1, entity-level scoring for NER), human evaluation protocols,
LLM-as-judge paradigms, the modern benchmark landscape, and decoding strategies
(greedy, beam search, temperature sampling, top-$k$, nucleus sampling, constrained
decoding).  Understanding evaluation deeply signals maturity---interviewers test whether
candidates know not just \emph{what} to measure, but \emph{why} certain metrics fail and
\emph{when} to distrust them.  A candidate who can derive BLEU, explain why ROUGE
favors extractive summaries, and articulate the biases of LLM-as-judge evaluation
demonstrates the kind of critical thinking that earns ``strong hire'' signals.
\end{tldr}

\bigskip

% ============================================================================
% SECTION 1: LANGUAGE MODELING METRICS
% ============================================================================
\section{Language Modeling Metrics}
\label{sec:lm_metrics}

The most fundamental evaluation question in NLP is: \emph{how well does the model
predict the next token?}  Language modeling metrics quantify this directly.

\subsection{Cross-Entropy Loss}

Given a sequence of $N$ tokens $w_1, w_2, \ldots, w_N$ and a model that assigns
probability $P_\theta(w_i \mid w_{<i})$ to each token conditioned on its left
context, the \textbf{cross-entropy loss} is:
\[
    H(\theta) = -\frac{1}{N} \sum_{i=1}^{N} \log_2 P_\theta(w_i \mid w_{<i})
\]

This measures the average number of \emph{bits} the model needs to encode each
token.  A perfect model that assigns probability 1 to every correct next token
achieves $H = 0$.  In practice, training minimizes cross-entropy using natural
logarithms ($\ln$), which differs from the information-theoretic convention
($\log_2$) by a constant factor: $H_{\text{nats}} = H_{\text{bits}} \times \ln 2$.

\subsection{Perplexity}

Perplexity exponentiates the cross-entropy to give a more interpretable number:

\begin{mathresult}{Perplexity}
\[
    \text{PP}(W) = \exp\!\left(-\frac{1}{N}\sum_{i=1}^{N}\ln P_\theta(w_i \mid w_{<i})\right)
    = 2^{H(\theta)}
\]
where $H(\theta)$ is the cross-entropy in bits.  Equivalently, using the chain rule
of probability:
\[
    \text{PP}(W) = P_\theta(w_1, w_2, \ldots, w_N)^{-1/N}
    = \left(\prod_{i=1}^{N} P_\theta(w_i \mid w_{<i})\right)^{-1/N}
\]
\end{mathresult}

\paragraph{Intuition: weighted average branching factor.}
Perplexity represents the effective number of equally likely choices the model
considers at each step.  A perplexity of 50 means the model is, on average, as
uncertain as if choosing uniformly among 50 options per token.  For English text:
\begin{itemize}
    \item A uniform distribution over a 50,000-token vocabulary gives
          $\text{PP} = 50{,}000$.
    \item A good n-gram model achieves $\text{PP} \approx 80\text{--}200$.
    \item GPT-2 (1.5B) achieved $\text{PP} \approx 18$ on WikiText-103.
    \item Modern LLMs achieve $\text{PP} < 10$ on standard benchmarks.
\end{itemize}
Lower is better, but the absolute number is only meaningful when comparing models
on the \emph{same} data with the \emph{same} tokenizer.

\subsection{Bits-Per-Character and Bits-Per-Byte}
\label{sec:bpc}

Because different tokenizers produce different numbers of tokens for the same
text, perplexity is not directly comparable across models with different
vocabularies.  \textbf{Bits-per-character} (BPC) and \textbf{bits-per-byte}
(BPB) solve this by normalizing to a tokenizer-independent unit:
\[
    \text{BPC} = \frac{-\sum_{i=1}^{N}\log_2 P_\theta(w_i \mid w_{<i})}{\text{total characters in text}}, \qquad
    \text{BPB} = \frac{-\sum_{i=1}^{N}\log_2 P_\theta(w_i \mid w_{<i})}{\text{total bytes (UTF-8)}}
\]
This enables fair comparison: a character-level model and a subword-level model
can be compared on the same BPC/BPB scale.

\begin{keyinsight}[Perplexity: Necessary but Not Sufficient]
Perplexity measures a model's ability to \emph{predict} text---not to
\emph{generate} useful, faithful, safe, or coherent text.  A model with
excellent perplexity might still hallucinate facts, produce toxic content, or
fail at instruction following.  This is precisely why the field moved to
downstream benchmarks (MMLU, HumanEval, GSM8K) and human evaluation as the
primary assessment tools for LLMs, while perplexity remains the dominant metric
for \emph{pretraining} optimization and scaling law research.
\end{keyinsight}

\subsection{Limitations of Perplexity}

\begin{enumerate}[label=\textbf{\arabic*.}]
    \item \textbf{Vocabulary dependence.}  Models with different tokenizers
          produce incomparable perplexities.  A BPE tokenizer with 32K vocabulary
          and one with 100K vocabulary will report different perplexities on
          identical text because they predict different token sequences.

    \item \textbf{No quality signal.}  Perplexity does not measure fluency,
          coherence, factuality, helpfulness, or safety.  Two models with
          identical perplexity can differ drastically in generation quality.

    \item \textbf{Domain sensitivity.}  Perplexity on Wikipedia does not predict
          performance on clinical notes, legal documents, or code.  Always report
          perplexity on a domain-matched held-out set.

    \item \textbf{Equal token weighting.}  Every token contributes equally, but
          predicting ``the'' is trivially easy while predicting a rare entity
          name is what matters for downstream tasks.

    \item \textbf{Not calibrated to utility.}  Improvements from $\text{PP}=15$ to
          $\text{PP}=14$ may matter more or less than from $\text{PP}=50$ to
          $\text{PP}=40$, depending on the downstream task.
\end{enumerate}

% ============================================================================
% SECTION 2: MACHINE TRANSLATION METRICS
% ============================================================================
\section{Machine Translation Metrics}
\label{sec:mt_metrics}

Machine translation evaluation has driven some of the most rigorous metric
development in NLP, because translation quality is multidimensional (adequacy,
fluency, terminology) and purely human evaluation is prohibitively expensive at
scale.

\subsection{BLEU: Bilingual Evaluation Understudy}
\label{sec:bleu}

BLEU (Papineni et al., ACL 2002) remains the most widely reported MT metric,
despite well-known limitations.  Every NLP interview candidate must be able to
derive it from scratch.

\paragraph{Step 1: Modified n-gram precision.}
For each n-gram order $n \in \{1,2,3,4\}$, compute the \emph{modified precision}
$p_n$, which clips the count of each n-gram in the candidate by its maximum count
in any reference translation:
\[
    p_n = \frac{\sum_{\text{n-gram} \in C}\min\!\bigl(\text{Count}_C(\text{n-gram}),\;\max_{R}\text{Count}_R(\text{n-gram})\bigr)}{\sum_{\text{n-gram} \in C}\text{Count}_C(\text{n-gram})}
\]
where $C$ is the candidate and $R$ iterates over reference translations.  The
clipping prevents a degenerate candidate like ``the the the the'' from achieving
perfect unigram precision against a reference containing ``the.''

\paragraph{Step 2: Geometric mean of n-gram precisions.}
BLEU combines precisions across n-gram orders using a weighted geometric mean:
\[
    \text{geo} = \exp\!\left(\sum_{n=1}^{4} w_n \ln p_n\right)
\]
where typically $w_n = \frac{1}{4}$ for all $n$ (uniform weighting).  Using the
geometric mean ensures that a zero precision at \emph{any} order drives the
entire score to zero, which penalizes candidates that match individual words but
fail to capture any longer phrases.

\paragraph{Step 3: Brevity penalty.}
Short candidates can trivially achieve high precision by generating only
high-confidence words.  The \textbf{brevity penalty} (BP) penalizes candidates
shorter than the reference:
\[
    \text{BP} = \begin{cases}
        1 & \text{if } |C| \geq |R| \\[4pt]
        \exp\!\left(1 - \frac{|R|}{|C|}\right) & \text{if } |C| < |R|
    \end{cases}
\]
where $|C|$ and $|R|$ are the candidate and (closest) reference lengths.

\begin{mathresult}{BLEU Score}
\[
    \text{BLEU} = \text{BP} \times \exp\!\left(\sum_{n=1}^{4} w_n \ln p_n\right)
\]
where $p_n$ is the modified $n$-gram precision, $w_n = \tfrac{1}{4}$ by default,
and $\text{BP} = \exp\!\bigl(\min(0,\; 1 - |R|/|C|)\bigr)$.  BLEU ranges from
0 to 1, though scores above 0.40 are considered very good for MT.
\end{mathresult}

\paragraph{Corpus-level vs.\ sentence-level.}
BLEU is designed as a \emph{corpus-level} metric: $p_n$ is computed by summing
counts across all sentence pairs before dividing.  Sentence-level BLEU is
problematic because short sentences often have $p_n = 0$ for higher $n$, which
drives the geometric mean to zero.  Smoothed variants (BLEU+1, SacreBLEU) add
small epsilon counts to handle this.

\subsection{Limitations of BLEU}

\begin{warningbox}
BLEU has serious known limitations that interviewers frequently test:
\begin{enumerate}[label=(\arabic*)]
    \item \textbf{Precision-only.}  BLEU measures precision (how many candidate
          n-grams appear in the reference) but not recall (how many reference
          n-grams appear in the candidate).  The brevity penalty is a crude
          recall proxy.
    \item \textbf{No semantic awareness.}  ``The cat sat on the mat'' and
          ``A feline rested upon the rug'' share few n-grams despite being
          near-paraphrases.  BLEU penalizes valid lexical variation.
    \item \textbf{Word-order insensitivity.}  BLEU-1 treats bags of words; even
          BLEU-4 only captures local 4-gram order, not global sentence structure.
    \item \textbf{Reference dependence.}  Quality depends heavily on reference
          count and quality.  A single reference captures only one valid
          translation among many.
    \item \textbf{Poor correlation at segment level.}  Individual sentence-level
          BLEU scores correlate poorly with human judgments.
    \item \textbf{Language bias.}  BLEU was designed for English and works poorly
          for agglutinative or morphologically rich languages where surface form
          variation is high.
\end{enumerate}
\end{warningbox}

\subsection{Beyond BLEU: Better MT Metrics}

\paragraph{METEOR} (Banerjee \& Lavie, 2005) addresses several BLEU limitations
by computing an alignment between candidate and reference tokens using exact
match, stemming, synonyms (via WordNet), and paraphrase tables.  It computes
both precision and recall, combines them with a weighted harmonic mean (favoring
recall), and applies a fragmentation penalty for non-contiguous matches.  METEOR
correlates better with human judgments than BLEU, particularly at the sentence
level.

\paragraph{chrF and chrF++} (Popovi\'{c}, 2015) compute character n-gram F-score
rather than word n-grams.  Because character n-grams are language-agnostic and
naturally capture morphological similarity, chrF works well across typologically
diverse languages without requiring stemming or synonym resources.  chrF++ adds
word unigram and bigram precision to the character-level metric.

\paragraph{COMET} (Rei et al., 2020) is a \emph{neural metric} that encodes the
source, candidate, and reference using a pretrained multilingual encoder (XLM-R)
and predicts a quality score using a regression head trained on human judgment
data (Direct Assessment scores from WMT).  COMET achieves substantially higher
correlation with human judgments than BLEU or METEOR and has become the de facto
standard in MT research since 2020.

\begin{keyinsight}[The Metric Evolution in MT]
The progression from BLEU to COMET mirrors the broader NLP evolution: from
surface-form heuristics (n-gram overlap) to learned, neural representations
trained on human judgments.  In interviews, demonstrating awareness of this
progression shows that you understand both the classical metrics (which you must
know) and the state of the art (which you should prefer in practice).
\end{keyinsight}

\subsection{MT Metric Comparison}

\begin{table}[H]
\centering
\caption{Comparison of machine translation metrics.}
\label{tab:mt_metrics}
\begin{tabularx}{\textwidth}{@{}l c c c X@{}}
\toprule
\textbf{Metric} & \textbf{Type} & \textbf{Human Corr.} & \textbf{Lang.\ Agnostic} & \textbf{Key Mechanism} \\
\midrule
BLEU       & Surface    & Moderate  & No  & Modified n-gram precision + brevity penalty \\
METEOR     & Surface+   & Good      & No  & Alignment with stems/synonyms, F-score \\
chrF       & Surface    & Good      & Yes & Character n-gram F-score \\
TER        & Surface    & Moderate  & Yes & Edit distance (insert/delete/substitute/shift) \\
BERTScore  & Neural     & High      & Yes & Token-level cosine similarity with BERT \\
COMET      & Neural     & Very High & Yes & Regression on source + candidate + reference embeddings \\
\bottomrule
\end{tabularx}
\end{table}

% ============================================================================
% SECTION 3: SUMMARIZATION METRICS
% ============================================================================
\section{Summarization Metrics}
\label{sec:summarization_metrics}

Summarization evaluation is notoriously difficult because a good summary must be
simultaneously fluent, coherent, relevant, concise, and---most critically---\emph{faithful} to the source document.

\subsection{ROUGE: Recall-Oriented Understudy for Gisting Evaluation}

ROUGE (Lin, 2004) is the dominant automatic metric for summarization.  Unlike
BLEU (which is precision-oriented), ROUGE is \emph{recall-oriented}: it measures
how much of the \emph{reference} summary is captured by the candidate.

\begin{mathresult}{ROUGE Variants}
\textbf{ROUGE-N} (unigram $N{=}1$, bigram $N{=}2$):
\[
    \text{ROUGE-N} = \frac{\sum_{\text{n-gram} \in \text{Ref}} \min\!\bigl(\text{Count}_{\text{Cand}}(\text{n-gram}),\;\text{Count}_{\text{Ref}}(\text{n-gram})\bigr)}{\sum_{\text{n-gram} \in \text{Ref}} \text{Count}_{\text{Ref}}(\text{n-gram})}
\]
\textbf{ROUGE-L} (longest common subsequence):
\[
    R_{\text{lcs}} = \frac{\text{LCS}(\text{Ref}, \text{Cand})}{|\text{Ref}|}, \quad
    P_{\text{lcs}} = \frac{\text{LCS}(\text{Ref}, \text{Cand})}{|\text{Cand}|}, \quad
    \text{ROUGE-L} = \frac{(1+\beta^2) R_{\text{lcs}} P_{\text{lcs}}}{R_{\text{lcs}} + \beta^2 P_{\text{lcs}}}
\]
where $\beta$ is typically set large (e.g., 1.2) to weight recall more heavily.
LCS captures the longest in-order (not necessarily contiguous) subsequence match.
\end{mathresult}

\paragraph{ROUGE-1 vs.\ ROUGE-2 vs.\ ROUGE-L.}
\begin{itemize}
    \item \textbf{ROUGE-1} measures unigram recall---content coverage at the word
          level.  It is the most commonly reported variant for summarization.
    \item \textbf{ROUGE-2} measures bigram recall---captures some local fluency
          and phrase-level overlap.
    \item \textbf{ROUGE-L} uses the longest common subsequence, which captures
          sentence-level structure without requiring contiguous n-gram matches.
\end{itemize}

\paragraph{Why ROUGE is recall-oriented (and why this matters).}
Summarization references are typically \emph{short}.  A system that copies the
entire source document would have perfect recall but produce a useless summary.
ROUGE counters this partly by using F-score variants that balance precision and
recall, but the canonical metric emphasizes recall.  The critical insight:
\emph{BLEU and ROUGE are complementary}---BLEU penalizes candidates that add
extraneous content (precision), while ROUGE penalizes candidates that miss
reference content (recall).

\subsection{BERTScore}

BERTScore (Zhang et al., ICLR 2020) moves beyond surface n-gram overlap by
computing similarity in \emph{embedding space}:

\begin{enumerate}
    \item Encode both candidate and reference tokens using a pretrained model
          (typically RoBERTa or DeBERTa), producing contextual embeddings.
    \item For each candidate token, find the maximum cosine similarity with any
          reference token (greedy matching):
          \[
              P_{\text{BERT}} = \frac{1}{|C|}\sum_{c_i \in C}\max_{r_j \in R}\cos(\vx_{c_i}, \vx_{r_j}), \quad
              R_{\text{BERT}} = \frac{1}{|R|}\sum_{r_j \in R}\max_{c_i \in C}\cos(\vx_{r_j}, \vx_{c_i})
          \]
    \item Compute the F1 score: $F_{\text{BERT}} = 2 P_{\text{BERT}} R_{\text{BERT}} / (P_{\text{BERT}} + R_{\text{BERT}})$.
\end{enumerate}

BERTScore captures paraphrases (``purchase'' $\leftrightarrow$ ``buy'') and
semantic equivalence that surface metrics miss.  It correlates better with human
judgments than ROUGE for abstractive summarization.

\subsection{Faithfulness Metrics}
\label{sec:faithfulness}

The most critical problem in summarization evaluation is \textbf{faithfulness}:
does the summary contain only information that is supported by the source
document?  Neither ROUGE nor BERTScore measure this---a summary that adds
plausible but fabricated facts can score well on both.

\begin{itemize}
    \item \textbf{FactCC} (Kr\'{y}sci\'{n}ski et al., 2020): A BERT-based binary
          classifier trained on synthetic claim-document pairs.  Classifies
          each summary sentence as consistent or inconsistent with the source.

    \item \textbf{SummaC} (Laban et al., 2022): Splits summaries into sentences
          and computes NLI (natural language inference) scores between each
          summary sentence and source sentences, aggregating into a consistency
          score.

    \item \textbf{QuestEval / QAFactEval}: Generates questions from the summary,
          answers them from the source, and checks agreement.  If the source
          cannot answer a question implied by the summary, that content is
          likely hallucinated.

    \item \textbf{LLM-based checking}: Prompt an LLM to verify whether each
          claim in the summary is supported by the source.  Scalable but
          inherits LLM biases and errors.
\end{itemize}

\subsection{Limitations of Summarization Metrics}

\begin{enumerate}[label=\textbf{\arabic*.}]
    \item \textbf{ROUGE favors extractive output.}  Because ROUGE measures
          n-gram overlap, systems that copy phrases verbatim from the source
          tend to score higher than abstractive systems that paraphrase.  This
          creates a perverse incentive toward less fluent but higher-overlap
          summaries.

    \item \textbf{No faithfulness signal.}  ROUGE and BERTScore reward overlap
          with the reference---not consistency with the source.  A hallucinated
          sentence that happens to overlap with the reference will be rewarded.

    \item \textbf{Reference quality bottleneck.}  Summarization references are
          often written by a single annotator.  Different valid summaries of the
          same document may have low ROUGE overlap with each other.

    \item \textbf{Length gaming.}  Longer summaries tend to have higher ROUGE
          recall.  Without length normalization, systems are incentivized to
          produce longer output.
\end{enumerate}

% ============================================================================
% SECTION 4: CLASSIFICATION METRICS FOR NLP
% ============================================================================
\section{Classification Metrics for NLP}
\label{sec:classification_metrics}

Classification is the backbone of many production NLP systems: spam detection,
sentiment analysis, intent classification, content moderation, and named entity
recognition all reduce to classification at some level.  The choice of metric
must match the task structure and business objective.

\subsection{Fundamental Metrics}

For a binary classifier with true positives (TP), false positives (FP), false
negatives (FN), and true negatives (TN):

\[
    \text{Precision} = \frac{\text{TP}}{\text{TP} + \text{FP}}, \qquad
    \text{Recall} = \frac{\text{TP}}{\text{TP} + \text{FN}}, \qquad
    F_1 = \frac{2 \cdot \text{Precision} \cdot \text{Recall}}{\text{Precision} + \text{Recall}}
\]

\paragraph{Precision vs.\ recall tradeoff.}
The key question is: \emph{what is more costly---a false positive or a false
negative?}  In spam filtering, a false positive (legitimate email in spam) is
worse than a false negative (spam in inbox), so precision matters more.  In
medical screening, a false negative (missed disease) is worse, so recall matters
more.  F1 balances both but treats them as equally important.

\subsection{Multi-class Averaging}

For $K$ classes, there are three standard averaging schemes:

\begin{itemize}
    \item \textbf{Macro F1:} Compute F1 per class, then average.  Treats all
          classes equally regardless of frequency.
          \[
              F_1^{\text{macro}} = \frac{1}{K}\sum_{k=1}^{K} F_1^{(k)}
          \]
    \item \textbf{Micro F1:} Aggregate TP, FP, FN across all classes, then
          compute F1 on the totals.  Equivalent to accuracy for single-label
          classification.
          \[
              F_1^{\text{micro}} = \frac{2\sum_k\text{TP}_k}{2\sum_k\text{TP}_k + \sum_k\text{FP}_k + \sum_k\text{FN}_k}
          \]
    \item \textbf{Weighted F1:} Compute F1 per class, average weighted by class
          frequency (support).  Accounts for class imbalance while still
          reflecting per-class performance.
\end{itemize}

\begin{warningbox}
\textbf{Micro vs.\ macro F1 in imbalanced NER datasets:}
In NER, entity types are highly imbalanced (e.g., ``O'' tags dominate).  Micro
F1 is dominated by the majority class and can be misleadingly high.  Macro F1
gives equal weight to rare entity types, revealing poor performance on minority
classes that micro F1 hides.  \textbf{Always report both}, and understand which
one the business objective aligns with.  If identifying rare entities (e.g.,
adverse drug reactions in medical text) is the primary goal, macro F1 is the
right metric.  If overall system throughput matters, micro F1 is appropriate.
\end{warningbox}

\subsection{Entity-Level F1 for NER}
\label{sec:ner_eval}

NER evaluation has a subtle but critical distinction: \textbf{token-level} vs.\
\textbf{entity-level} evaluation.

\paragraph{Token-level F1.}
Evaluate each token's predicted BIO tag independently.  A prediction of B-PER
followed by I-PER for ``Barack Obama'' scores two correct tokens.  Problem:
partial matches inflate scores.  If the model predicts ``Barack'' as B-PER but
misses ``Obama,'' token-level scoring gives 50\% recall for that entity, which
is misleading---the entity was \emph{not} correctly extracted.

\paragraph{Entity-level F1 (strict matching).}
An entity is counted as a true positive only if \textbf{both the boundary and
the type are exactly correct}:
\begin{itemize}
    \item Predicted span must match the gold span exactly (same start and end
          token).
    \item Predicted entity type must match the gold entity type.
    \item Any mismatch in boundary or type counts as both a false positive
          (for the predicted entity) and a false negative (for the gold entity).
\end{itemize}

This is the standard evaluation in CoNLL-2003 and most NER benchmarks.  Entity-level
F1 is always lower than token-level F1 and gives a more honest picture of system
quality.

\paragraph{Partial matching (SemEval scheme).}
For tasks where partial credit is reasonable (e.g., extracting product names
where ``iPhone 15 Pro'' partially matching ``iPhone 15 Pro Max'' has value), the
SemEval evaluation scheme defines:
\begin{itemize}
    \item \textbf{Exact}: boundary and type both correct.
    \item \textbf{Partial}: overlapping boundary, correct type (half credit).
    \item \textbf{Type only}: correct type regardless of boundary.
    \item \textbf{Strict}: same as exact (no partial credit).
\end{itemize}

\subsection{Multi-Label Classification Metrics}

When each instance can have multiple labels (e.g., topic tagging, multi-label
sentiment), standard metrics extend as follows:
\begin{itemize}
    \item \textbf{Hamming loss:} Fraction of labels that are incorrectly
          predicted (either false positive or false negative), averaged over all
          instances and labels.
    \item \textbf{Subset accuracy (exact match ratio):} Fraction of instances
          where the \emph{entire} predicted label set exactly matches the gold
          label set.  Very strict---a single missed or extra label counts as
          wrong.
    \item \textbf{Label-wise F1:} Compute F1 per label, then macro/micro
          average.
\end{itemize}

% ============================================================================
% SECTION 5: HUMAN EVALUATION
% ============================================================================
\section{Human Evaluation}
\label{sec:human_eval}

Automatic metrics are proxies.  Human evaluation is the gold standard---and also
the most expensive, slowest, and hardest-to-reproduce evaluation method.
Understanding when and how to deploy it is a mark of evaluation maturity.

\subsection{Evaluation Dimensions}

Human evaluators typically rate generated text along multiple orthogonal
dimensions:

\begin{itemize}
    \item \textbf{Fluency:} Is the text grammatically correct and natural?
    \item \textbf{Coherence:} Does the text flow logically from sentence to
          sentence?
    \item \textbf{Relevance:} Does the output address the input query/task?
    \item \textbf{Faithfulness / Factuality:} Is the content consistent with the
          source material (summarization, RAG) or factually correct (open-ended
          generation)?
    \item \textbf{Harmlessness:} Does the output avoid toxic, biased, or
          dangerous content?
    \item \textbf{Helpfulness:} Does the output actually help the user accomplish
          their goal?
\end{itemize}

\subsection{Evaluation Protocols}

\paragraph{Likert Scale Rating.}
Evaluators rate each output on a numerical scale (e.g., 1--5).  Simple but
suffers from annotator calibration differences: one annotator's ``3'' may be
another's ``4.''  Anchoring the scale with examples (``1 = incomprehensible,
5 = publication quality'') helps but does not eliminate subjectivity.

\paragraph{Pairwise Comparison.}
Evaluators are shown two outputs (from different systems) for the same input and
choose which is better (or ``tie'').  More reliable than absolute ratings because
it removes calibration differences---humans are better at relative judgments than
absolute scores.  This is the protocol used by Chatbot Arena.

\paragraph{Ranking.}
Evaluators rank $k$ outputs from best to worst.  More informative than pairwise
but quadratically more comparisons to obtain a full ranking.  Best-worst scaling
(BWS) is a popular compromise: evaluators select only the best and worst from
each group.

\subsection{Inter-Annotator Agreement}

Because human evaluation is subjective, measuring agreement between annotators
is critical:
\begin{itemize}
    \item \textbf{Cohen's kappa} (for 2 annotators): $\kappa = (p_o - p_e) / (1 - p_e)$,
          where $p_o$ is observed agreement and $p_e$ is chance agreement.
          $\kappa > 0.8$ is excellent; $\kappa < 0.4$ is poor.
    \item \textbf{Krippendorff's alpha}: Generalizes to multiple annotators,
          multiple categories, and handles missing data.  The recommended
          standard for NLP annotation studies.
    \item \textbf{Fleiss' kappa}: Extends Cohen's kappa to more than 2
          annotators when each item is rated by the same number of annotators.
\end{itemize}

Low agreement does not necessarily mean the annotations are bad---it may mean
the task is genuinely ambiguous or the guidelines are underspecified.

\subsection{Crowdsourced Evaluation at Scale}

\paragraph{LMSYS Chatbot Arena.}
The most influential LLM evaluation platform (since 2023).  Users submit
prompts, receive anonymous responses from two models, and vote for the better
one.  Votes are aggregated into an Elo rating system (adapted from chess).  Key
strengths: real user prompts, massive scale (millions of votes), no prompt
engineering by model creators.  Key limitations: population bias (tech-savvy
English speakers), prompt distribution does not match all use cases, Elo
sensitivity to prompt difficulty distribution.

\paragraph{MT-Bench.}
A curated set of 80 multi-turn questions spanning 8 categories (writing,
reasoning, math, coding, extraction, STEM, humanities, roleplay).  Responses are
scored 1--10 by GPT-4 as judge.  Useful for rapid evaluation during development
but limited by prompt set size and GPT-4 judge biases.

% ============================================================================
% SECTION 6: LLM-AS-JUDGE
% ============================================================================
\section{LLM-as-Judge}
\label{sec:llm_judge}

Using a powerful LLM (GPT-4, Claude) to evaluate the outputs of other models has
become a dominant evaluation paradigm because it scales better than human
evaluation while correlating reasonably well with human preferences.

\subsection{Evaluation Modes}

\begin{itemize}
    \item \textbf{Pointwise scoring:} The judge LLM rates a single response on a
          scale (e.g., 1--10) given a rubric.  Simple but suffers from
          calibration drift and positional sensitivity within the prompt.
    \item \textbf{Pairwise comparison:} The judge compares two responses and
          selects the better one.  More robust than pointwise because relative
          judgments are easier.  This is the protocol used in AlpacaEval.
    \item \textbf{Listwise ranking:} The judge ranks multiple responses.  Most
          informative but most prone to position bias (the judge may favor
          responses presented first or last).
    \item \textbf{Reference-guided:} The judge is given a reference answer and
          evaluates how well the candidate matches it.  Useful for factual tasks.
\end{itemize}

\subsection{Known Biases}

\begin{warningbox}
LLM-as-judge evaluation has systematic biases that can silently corrupt evaluation
results:
\begin{enumerate}[label=(\arabic*)]
    \item \textbf{Verbosity bias:} LLM judges tend to prefer longer, more detailed
          responses, even when brevity is more appropriate.  Models that pad
          responses with unnecessary elaboration are systematically overrated.
    \item \textbf{Position bias:} When comparing two responses, LLM judges may
          prefer the response presented first (primacy) or second (recency),
          regardless of quality.  Mitigation: evaluate both orderings and average.
    \item \textbf{Self-preference bias:} GPT-4 as judge tends to prefer
          GPT-4-generated responses over other models' outputs, even when human
          evaluators disagree.  Similarly for other model families.
    \item \textbf{Style over substance:} LLM judges may reward well-formatted,
          confident-sounding responses over factually correct but less polished
          ones.
    \item \textbf{Sycophancy in evaluation:} Some judges agree too readily with
          the prompt framing, inflating scores when the prompt implies the
          response should be good.
\end{enumerate}
\end{warningbox}

\subsection{Calibration and Best Practices}

\begin{itemize}
    \item \textbf{Ground against human preferences:} Always validate that the
          LLM judge agrees with human annotators on a calibration set before
          trusting it on new data.  Report the agreement rate.
    \item \textbf{Swap position:} For pairwise evaluation, run each comparison
          in both orders (A-B and B-A) and flag disagreements.
    \item \textbf{Use rubrics:} Detailed scoring rubrics reduce variance.
          ``Rate helpfulness 1--10'' is worse than a rubric with examples for
          each score level.
    \item \textbf{Multiple judges:} Use multiple LLM judges (e.g., GPT-4 +
          Claude + Gemini) and check agreement, analogous to inter-annotator
          agreement for human evaluation.
    \item \textbf{Report the judge:} Always state which LLM and version was used
          as judge---evaluation results are not reproducible without this.
\end{itemize}

\begin{keyinsight}[LLM-as-Judge: Useful but Not Authoritative]
LLM-as-judge is best understood as a \emph{noisy, scalable proxy} for human
evaluation.  It is invaluable during development for rapid iteration (catching
major regressions, A/B testing prompt variants) but should not be the sole
basis for high-stakes decisions like model deployment or benchmark claims.  For
any result you plan to publish or deploy, validate with human evaluation on
a representative sample.
\end{keyinsight}

% ============================================================================
% SECTION 7: BENCHMARK LANDSCAPE
% ============================================================================
\section{The Benchmark Landscape}
\label{sec:benchmarks}

Benchmarks serve as standardized tests for comparing models, but they are
imperfect proxies for real-world utility.  A mature candidate knows what each
benchmark tests, why it exists, and where it fails.

\subsection{Major LLM Benchmarks}

\begin{table}[H]
\centering
\caption{Major benchmarks for LLM evaluation.}
\label{tab:benchmarks}
\begin{tabularx}{\textwidth}{@{}l l l X@{}}
\toprule
\textbf{Benchmark} & \textbf{Format} & \textbf{Size} & \textbf{What It Tests} \\
\midrule
MMLU         & 4-way MC   & 15K Qs, 57 subjects & Broad knowledge (STEM, humanities, social science) \\
HumanEval    & Code gen.  & 164 problems         & Python function completion, test-case passing \\
MBPP         & Code gen.  & 974 problems          & Basic Python programming \\
GSM8K        & Free-form  & 8.5K problems         & Grade-school math reasoning (multi-step) \\
MATH         & Free-form  & 12.5K problems        & Competition-level mathematics \\
ARC          & 4-way MC   & 7.7K questions        & Grade-school / challenge science reasoning \\
HellaSwag    & 4-way MC   & 10K questions         & Commonsense sentence completion \\
WinoGrande   & Binary     & 44K problems          & Commonsense coreference resolution \\
TruthfulQA   & MC / gen.  & 817 questions         & Resistance to common misconceptions \\
BigBench     & Various    & 200+ tasks            & Diverse capabilities (reasoning, translation, QA) \\
MT-Bench     & Open-ended & 80 multi-turn Qs      & Multi-turn conversation quality (LLM-judged) \\
Chatbot Arena & Pairwise  & Millions of votes     & Real-user preference (Elo rating) \\
\bottomrule
\end{tabularx}
\end{table}

\subsection{Benchmark Pitfalls}

\paragraph{Data contamination.}
LLMs are trained on vast web corpora that may contain benchmark questions and
answers.  A model that has ``seen'' GSM8K problems during pretraining will
overfit to them, inflating its reported score.  Contamination is difficult to
detect---even deduplication against known benchmarks cannot catch paraphrased
versions.  Mitigation strategies include: held-out test sets never published
online, dynamic benchmarks with new questions per evaluation round (LiveBench),
and contamination analysis (n-gram overlap between training data and benchmarks).

\paragraph{Benchmark saturation.}
As models improve, older benchmarks lose discriminative power.  HellaSwag and
WinoGrande are near ceiling for frontier models.  Even MMLU has become less
informative as top models exceed 90\%.  The field must continuously develop
harder benchmarks---MMLU-Pro, GPQA (PhD-level questions), and FrontierMath
represent this trend.

\paragraph{Gaming and Goodhart's law.}
``When a measure becomes a target, it ceases to be a good measure.''  Models can
be optimized specifically for benchmark performance (training on similar
distributions, prompt engineering) without corresponding real-world improvement.
Chatbot Arena partially mitigates this by using real, unpredictable user prompts.

\paragraph{Narrow coverage.}
No single benchmark captures the full range of capabilities that matter:
safety, instruction following, long-context reasoning, multilingual performance,
tool use, agentic planning, and domain-specific knowledge all require specialized
evaluation.

\begin{keyinsight}[Benchmarks Are Necessary but Unreliable]
Interviewers are not testing whether you memorized benchmark scores---they are
testing whether you understand the \emph{limitations} of benchmarks.  The
strongest answer to ``how would you evaluate your model?'' is never ``we scored
X on MMLU'' but rather a multi-layered evaluation strategy: automatic metrics
for development iteration, targeted benchmarks for capability assessment, human
evaluation for quality validation, and production metrics (engagement, task
completion rate) for real-world impact.
\end{keyinsight}

% ============================================================================
% SECTION 8: DECODING STRATEGIES
% ============================================================================
\section{Decoding Strategies}
\label{sec:decoding}

A language model defines a probability distribution over next tokens.
\emph{Decoding} is the process of selecting tokens from this distribution to
generate output text.  The choice of decoding strategy profoundly affects output
quality, diversity, and computational cost.

\subsection{Greedy Decoding}

At each step, select the token with the highest probability:
\[
    w_t = \argmax_{w} P_\theta(w \mid w_{<t})
\]
\textbf{Advantages:} Simplest possible strategy; fast (no branching).
\textbf{Disadvantages:} Produces repetitive, dull text because it always takes
the locally optimal choice, which often leads to loops (``the the the...'' or
repeating sentence patterns).  It also misses globally optimal sequences---the
best overall sentence may start with a less-probable first word.

\subsection{Beam Search}

Maintain $B$ candidate sequences (beams) at each step, expanding each by the
top-$k$ tokens and keeping the top $B$ overall:

\begin{enumerate}
    \item Initialize with the $B$ most probable first tokens.
    \item At each subsequent step, expand each beam by its top-$V$ continuations
          (often $V = B$), yielding up to $B \times V$ candidates.
    \item Score each candidate by cumulative log-probability and keep the top $B$.
    \item Terminate beams that produce an end-of-sequence token.
    \item Return the highest-scoring complete sequence.
\end{enumerate}

\paragraph{Length normalization.}
Raw log-probability is biased toward shorter sequences (fewer terms summed).
Length normalization divides by $|y|^\alpha$ where $\alpha \in [0.6, 1.0]$:
\[
    \text{score}(y) = \frac{\log P(y)}{|y|^\alpha}
\]
Setting $\alpha = 0$ ignores length; $\alpha = 1$ gives per-token
log-probability.  In practice, $\alpha = 0.6\text{--}0.8$ works well for
machine translation.

\paragraph{When to use beam search.}
Beam search excels for tasks with a clear ``correct'' answer---machine
translation, speech recognition, structured output generation.  It produces
reliable, high-likelihood output but lacks diversity.  For open-ended generation
(chatbots, creative writing), beam search produces bland, repetitive text.

\subsection{Temperature Sampling}

\emph{Sampling} introduces randomness by drawing from the distribution rather
than taking the argmax.  \textbf{Temperature} controls the sharpness of this
distribution:

\begin{mathresult}{Temperature Scaling}
Given logits $z_i$ for token $i$, temperature-scaled sampling draws from:
\[
    P(w_i) = \frac{\exp(z_i / T)}{\sum_j \exp(z_j / T)}
\]
\begin{itemize}
    \item $T = 1$: the model's original distribution (no modification).
    \item $T < 1$: sharper distribution (more confident, less diverse).
          As $T \to 0$, sampling becomes greedy decoding.
    \item $T > 1$: flatter distribution (more random, more diverse).
          As $T \to \infty$, sampling becomes uniform random over the vocabulary.
\end{itemize}
\end{mathresult}

\paragraph{Why temperature works.}
Consider logits $[2.0, 1.0, 0.5]$.  At $T{=}1$, probabilities are approximately
$[0.50, 0.27, 0.13]$.  At $T{=}0.5$, they become $[0.71, 0.21, 0.08]$---the
winner dominates more.  At $T{=}2.0$, they become $[0.39, 0.33, 0.28]$---nearly
uniform.  Temperature does not change the \emph{ranking} of tokens, only the
\emph{magnitude of differences} between their probabilities.

\subsection{Top-$k$ Sampling}

At each step, restrict sampling to the $k$ most probable tokens, redistributing
probability mass among them:

\[
    P'(w_i) = \begin{cases}
        P(w_i) \;/\; \sum_{j \in \text{top-}k} P(w_j) & \text{if } w_i \in \text{top-}k \\[4pt]
        0 & \text{otherwise}
    \end{cases}
\]

\textbf{Problem:}  A fixed $k$ is not adaptive.  When the model is confident
(one token has 95\% probability), $k{=}40$ still includes 39 unlikely tokens.
When the model is uncertain (flat distribution), $k{=}40$ may exclude plausible
tokens.

\subsection{Top-$p$ (Nucleus) Sampling}

Top-$p$ sampling (Holtzman et al., ICLR 2020) addresses the fixed-$k$ problem by
dynamically selecting the smallest set of tokens whose cumulative probability
exceeds a threshold $p$:

\[
    V_p = \min \{V' \subseteq \mathcal{V} : \sum_{w \in V'} P(w) \geq p\}
\]
where tokens are sorted by descending probability and added until the cumulative
sum reaches $p$.  Then sample from $V_p$ with renormalized probabilities.

\textbf{Why this is better than top-$k$:}  The size of $V_p$ adapts to the
model's uncertainty.  When the model is confident, $V_p$ might contain only
1--3 tokens.  When uncertain, $V_p$ might contain hundreds.  Typical values:
$p \in [0.9, 0.95]$.

\subsection{Repetition Penalties}

LLMs tend to repeat themselves, especially with greedy or low-temperature
decoding.  Several mechanisms combat this:

\begin{itemize}
    \item \textbf{Repetition penalty} (Keskar et al., 2019): Divide the logit
          of any previously generated token by a penalty factor $\theta > 1$:
          $z'_i = z_i / \theta$ if token $i$ has appeared before.  Simple and
          effective.
    \item \textbf{Frequency penalty:} Subtract a term proportional to how many
          times a token has appeared: $z'_i = z_i - \alpha \cdot \text{count}(i)$.
          Used in OpenAI's API.
    \item \textbf{Presence penalty:} Subtract a fixed amount if the token has
          appeared \emph{at all}: $z'_i = z_i - \beta \cdot \mathbbm{1}[\text{count}(i) > 0]$.
          Encourages topic diversity without penalizing appropriate repetition.
\end{itemize}

\subsection{Constrained Decoding}

For tasks requiring structured output (JSON, SQL, code, XML), \textbf{constrained
decoding} restricts the token selection to only valid continuations:

\begin{itemize}
    \item \textbf{Grammar-based (GBNF):} Define a formal grammar for the target
          format.  At each step, mask out tokens that would violate the grammar.
          Guarantees syntactically valid output.
    \item \textbf{Regex-based:} Similar idea using regular expressions to define
          valid token sequences.
    \item \textbf{JSON mode:} Many inference frameworks (vLLM, Outlines, llama.cpp)
          implement JSON schema-constrained decoding, ensuring every generated
          token is a valid continuation of a JSON document matching the schema.
    \item \textbf{Logit bias:} Manually boost or suppress specific tokens.
          Coarser than grammar-based methods but simpler to implement.
\end{itemize}

\subsection{Speculative Decoding}
\label{sec:speculative}

Standard autoregressive decoding is memory-bandwidth bound: generating one token
requires loading the entire model's weights from memory.
\textbf{Speculative decoding} (Leviathan et al., 2023; Chen et al., 2023)
addresses this with a draft-and-verify approach:

\begin{enumerate}
    \item A small, fast \emph{draft model} generates $k$ candidate tokens
          autoregressively.
    \item The large \emph{target model} evaluates all $k$ tokens in parallel
          (a single forward pass).
    \item Accept draft tokens from left to right as long as the target model
          agrees (using a rejection sampling scheme).
    \item On the first disagreement, resample from the target model's adjusted
          distribution.
\end{enumerate}

The critical property: speculative decoding produces \textbf{exactly the same
distribution} as standard decoding from the target model.  The speedup comes from
amortizing the target model's forward pass across multiple tokens.  Typical
speedups: 2--3$\times$ when draft acceptance rate is high.

\subsection{Decoding Strategy Comparison}

\begin{table}[H]
\centering
\caption{Comparison of decoding strategies.}
\label{tab:decoding}
\begin{tabularx}{\textwidth}{@{}l c c c X@{}}
\toprule
\textbf{Strategy} & \textbf{Quality} & \textbf{Diversity} & \textbf{Latency} & \textbf{Best Use Case} \\
\midrule
Greedy          & Moderate & None   & Lowest  & Debugging, deterministic tasks \\
Beam search     & High     & Low    & High    & Translation, speech recognition \\
Temperature     & Tunable  & Tunable & Low    & General generation (tune $T$) \\
Top-$k$         & Good     & Moderate & Low   & Creative writing (fixed cutoff) \\
Top-$p$ (nucleus) & Good   & Good   & Low    & Open-ended generation (adaptive) \\
Constrained     & Task-dep.& None   & Moderate & JSON, SQL, structured output \\
Speculative     & Same*    & Same*  & Lower  & Production serving (*matches target model) \\
\bottomrule
\end{tabularx}
\end{table}

\begin{keyinsight}[Choosing the Right Decoding Strategy]
The choice of decoding strategy should match the task:
\begin{itemize}
    \item \textbf{Deterministic tasks} (classification, extraction, factual QA):
          greedy or low temperature ($T \leq 0.3$).
    \item \textbf{Constrained generation} (translation, summarization): beam
          search with length normalization.
    \item \textbf{Open-ended generation} (chatbots, creative writing): nucleus
          sampling ($p = 0.9\text{--}0.95$) with moderate temperature
          ($T = 0.7\text{--}1.0$).
    \item \textbf{Structured output} (JSON, code): constrained decoding with
          grammar enforcement.
    \item \textbf{Production serving at scale}: add speculative decoding for
          throughput improvement.
\end{itemize}
A candidate who can articulate these distinctions---rather than defaulting to
``just set temperature to 0.7''---demonstrates production awareness.
\end{keyinsight}

% ============================================================================
% SECTION 9: INTERVIEW QUESTIONS
% ============================================================================
\section{Interview Questions}
\label{sec:eval_interview_questions}

\vspace{0.5em}

%% QUESTION 1 %%
\begin{interviewq}
\textbf{Q: What is BLEU? Derive it step by step. What are its limitations?}

\badgeLfive{L5} \quad \badgetype{Mathematical} \quad \badgefreq{\freqhigh}

\stronganswer
BLEU (Bilingual Evaluation Understudy) is an automatic metric for evaluating
machine translation quality by measuring n-gram precision between a candidate
translation and one or more reference translations.

\textbf{Derivation:}
\begin{enumerate}
    \item \textbf{Modified n-gram precision $p_n$:} For each n-gram in the
          candidate, count how many times it appears, but clip this count by
          the maximum number of times it appears in any single reference.
          This prevents gaming (e.g., repeating ``the'' 20 times).
          \[
              p_n = \frac{\sum_{\text{n-gram} \in C}\min(\text{Count}_C,\;\max_R \text{Count}_R)}{\sum_{\text{n-gram} \in C}\text{Count}_C}
          \]

    \item \textbf{Geometric mean:} Combine $p_1$ through $p_4$ using a weighted
          geometric mean with uniform weights $w_n = 1/4$:
          $\exp\!\bigl(\sum_{n=1}^{4}\frac{1}{4}\ln p_n\bigr)$.
          This ensures that failure at any n-gram order zeroes the entire score.

    \item \textbf{Brevity penalty:}
          $\text{BP} = \exp\!\bigl(\min(0, 1 - |R|/|C|)\bigr)$.
          Penalizes candidates shorter than the reference to prevent trivially
          high precision from short outputs.

    \item \textbf{Final score:}
          $\text{BLEU} = \text{BP} \times \exp\!\bigl(\sum_{n=1}^{4}w_n \ln p_n\bigr)$.
\end{enumerate}

\textbf{Limitations:}
(1)~Precision-only: no recall component beyond the crude brevity penalty.
(2)~No semantic awareness: valid paraphrases are penalized.
(3)~Poor segment-level correlation with human judgments.
(4)~Language-biased: works poorly for morphologically rich languages.
(5)~Insensitive to word order beyond short n-gram windows.
(6)~Reference-dependent: quality limited by reference count and quality.

Modern alternatives like COMET (neural, trained on human judgments) achieve
much higher correlation with human evaluation and should be preferred in
practice, but BLEU remains the lingua franca of MT reporting.

\redflags
\begin{itemize}
    \item Cannot derive the formula---just says ``it's n-gram overlap.''
    \item Does not mention the brevity penalty or explain why it exists.
    \item Does not mention the clipping in modified precision.
    \item Claims BLEU measures recall (it measures precision).
    \item Cannot name any limitation.
\end{itemize}

\interviewtest{Mathematical rigor, ability to derive metrics from scratch, and
critical awareness of metric limitations.  This question tests whether the
candidate understands metrics deeply enough to know when to trust them.}

\goodgreat
\begin{itemize}
    \item \badgeLfive{L5}: Derives all four components correctly, gives at least
          three limitations.
    \item \badgeLsix{L6}: Contrasts with ROUGE (recall-oriented), explains
          why BLEU is a corpus-level metric (sentence-level problems), discusses
          SacreBLEU for reproducibility, and mentions COMET as a superior
          alternative.
    \item \badgeLseven{L7}: Discusses the history of MT evaluation (BLEU's
          impact on driving MT progress despite its flaws), connects to the
          broader theme of Goodhart's law in metrics, and explains how WMT
          shared tasks transitioned from BLEU to neural metrics.
\end{itemize}

\followups{What is SacreBLEU and why was it created?  How does BLEU compare to
ROUGE?  What is COMET and why is it better?  How would you evaluate MT for a
morphologically rich language like Finnish?}
\end{interviewq}

\vspace{0.5em}

%% QUESTION 2 %%
\begin{interviewq}
\textbf{Q: What is perplexity? How does it relate to cross-entropy? What are its
limitations?}

\badgeLfive{L5} \quad \badgetype{Mathematical} \quad \badgefreq{\freqhigh}

\stronganswer
Perplexity is the standard evaluation metric for language models.  Given a
sequence $w_1, \ldots, w_N$ and a model assigning probability
$P_\theta(w_i \mid w_{<i})$, perplexity is:
\[
    \text{PP} = \exp\!\Bigl(-\frac{1}{N}\sum_{i=1}^{N}\ln P_\theta(w_i \mid w_{<i})\Bigr)
\]

\textbf{Connection to cross-entropy:}  The term inside the exponential is the
\emph{cross-entropy} $H$ (in nats).  Perplexity is simply $\text{PP} = e^H$
(or $2^{H_{\text{bits}}}$ in the information-theoretic convention using
$\log_2$).  Cross-entropy measures the average negative log-likelihood per
token; perplexity exponentiates it to give the ``effective vocabulary size'' the
model is choosing from at each step.  A perplexity of 20 means the model is, on
average, as uncertain as if choosing uniformly among 20 options.

\textbf{Connection to bits-per-character:}  To compare models with different
tokenizers, normalize by the number of characters or bytes rather than tokens.
This gives bits-per-character (BPC) or bits-per-byte (BPB), which are
tokenizer-independent.

\textbf{Limitations:}
\begin{enumerate}
    \item Not comparable across tokenizers (different vocabularies $\Rightarrow$
          different number of prediction events).
    \item Does not measure generation quality---a low-perplexity model can still
          hallucinate, produce toxic content, or fail at instruction following.
    \item Domain-sensitive: perplexity measured on one domain does not predict
          performance on another.
    \item All tokens weighted equally, but content words matter more than
          function words for most applications.
\end{enumerate}

For LLM evaluation, perplexity is useful for pretraining monitoring and scaling
law research, but downstream benchmarks and human evaluation are needed to assess
actual model utility.

\redflags
\begin{itemize}
    \item Cannot state the formula or the connection to cross-entropy.
    \item Says ``lower perplexity means better generation'' without qualification.
    \item Does not mention the vocabulary dependence issue.
    \item Confuses perplexity with accuracy.
\end{itemize}

\interviewtest{Mathematical understanding of language model evaluation, ability
to connect related concepts (cross-entropy, bits, information theory), and
awareness of metric limitations.}

\goodgreat
\begin{itemize}
    \item \badgeLfive{L5}: States the formula, gives the branching factor
          intuition, names two limitations.
    \item \badgeLsix{L6}: Derives the cross-entropy connection, discusses
          bits-per-byte as a tokenizer-independent alternative, and connects
          to downstream evaluation.
    \item \badgeLseven{L7}: Discusses how perplexity relates to scaling laws
          ($L = aC^{-\alpha}$ where $L$ is cross-entropy loss), whether
          low perplexity implies calibrated probabilities (not necessarily),
          and when perplexity improvements stop translating to downstream gains.
\end{itemize}

\followups{How would you compare perplexity between GPT-2 (BPE, 50K vocab) and
a character-level model?  What is bits-per-byte?  How does perplexity relate to
scaling laws?  At what point do perplexity improvements stop mattering for
downstream tasks?}
\end{interviewq}

\vspace{0.5em}

%% QUESTION 3 %%
\begin{interviewq}
\textbf{Q: How would you evaluate a summarization system? Specifically, how would
you measure faithfulness?}

\badgeLfive{L5} \quad \badgetype{System Design} \quad \badgefreq{\freqhigh}

\stronganswer
Summarization evaluation requires a \emph{multi-metric} approach because no
single metric captures all quality dimensions:

\textbf{1.\ Content coverage (automatic):}
ROUGE-1 (unigram recall), ROUGE-2 (bigram recall), and ROUGE-L (longest common
subsequence) measure how much reference content the summary captures.  BERTScore
adds semantic matching.  But these only measure overlap with a reference, not
consistency with the source.

\textbf{2.\ Faithfulness (the critical dimension):}
This is the most important and hardest dimension.  A summary that sounds fluent
but contains hallucinated facts is worse than useless.  I would measure
faithfulness with:
\begin{itemize}
    \item \textbf{NLI-based checking (SummaC):} Decompose the summary into
          sentences, run NLI between each summary sentence and the source.
          Sentences classified as ``contradiction'' flag hallucinations.
    \item \textbf{QA-based checking (QuestEval):} Generate questions from the
          summary, answer them using only the source, and verify agreement.
          Unanswerable questions indicate hallucinated content.
    \item \textbf{LLM-based verification:} Prompt a strong LLM to identify
          claims in the summary that are not supported by the source.  Scalable
          but imperfect.
\end{itemize}

\textbf{3.\ Fluency and coherence (human evaluation):}
Automatic metrics cannot reliably assess whether the summary reads naturally and
flows logically.  Human evaluation with Likert scales on fluency (1--5) and
coherence (1--5) is necessary for final quality assessment.

\textbf{4.\ Production metrics:}
In deployment, track user engagement (click-through rate on summaries), task
completion rate (did the summary help the user decide?), and flagged errors
(user reports of incorrect summaries).

The key architectural insight is that ROUGE measures overlap with a
\emph{reference} while faithfulness measures consistency with the \emph{source}
---these are different signals and both must be evaluated.

\redflags
\begin{itemize}
    \item Only mentions ROUGE with no discussion of faithfulness.
    \item Proposes BLEU for summarization (BLEU is precision-oriented; ROUGE's
          recall focus is more appropriate).
    \item Treats ``it looks good to me'' as a valid evaluation strategy.
    \item Cannot distinguish between reference overlap and source faithfulness.
\end{itemize}

\interviewtest{Multi-dimensional evaluation thinking, awareness of the
faithfulness problem, practical evaluation pipeline design.}

\goodgreat
\begin{itemize}
    \item \badgeLfive{L5}: Proposes ROUGE + human evaluation + at least one
          faithfulness method.
    \item \badgeLsix{L6}: Describes the full NLI/QA-based faithfulness pipeline,
          explains why ROUGE cannot detect hallucination, and proposes production
          monitoring metrics.
    \item \badgeLseven{L7}: Designs a complete evaluation framework including
          automated regression testing (ROUGE threshold gates), faithfulness
          auditing (sampled human annotation), production monitoring (error
          rates, user feedback loops), and discusses the cost-quality tradeoff
          in evaluation.
\end{itemize}

\followups{What is the difference between \emph{extrinsic} and \emph{intrinsic}
hallucination?  How would you build a faithfulness classifier?  What are the
limitations of NLI-based faithfulness checking?}
\end{interviewq}

\vspace{0.5em}

%% QUESTION 4 %%
\begin{interviewq}
\textbf{Q: What are the problems with LLM-as-judge evaluation? When would you
use it and when would you not?}

\badgeLfive{L5} \quad \badgetype{Conceptual} \quad \badgefreq{\freqmed}

\stronganswer
LLM-as-judge (using a powerful LLM like GPT-4 to evaluate other models' outputs)
has become popular because it scales much better than human evaluation while
maintaining reasonable correlation with human preferences.  However, it has
systematic biases:

\textbf{Known biases:}
\begin{enumerate}
    \item \textbf{Verbosity bias:} LLM judges prefer longer, more detailed
          responses even when conciseness is appropriate.  This systematically
          inflates scores for verbose models.
    \item \textbf{Position bias:} In pairwise comparison, the judge may prefer
          whichever response appears first (or second), regardless of quality.
    \item \textbf{Self-preference:} GPT-4 tends to prefer GPT-4-generated text;
          Claude prefers Claude-generated text.  The judge's family bias
          contaminates the evaluation.
    \item \textbf{Style over substance:} Well-formatted, confident responses are
          favored over correct but less polished ones.
\end{enumerate}

\textbf{When to use it:}  Rapid development iteration (comparing prompt
variants, catching regressions), evaluating subjective dimensions (helpfulness,
tone) where automated metrics are poor, and as a first-pass filter before
expensive human evaluation.

\textbf{When NOT to use it:}  Final deployment decisions, published benchmark
claims, evaluating factual accuracy (the judge may share the same factual errors
as the candidate), safety-critical applications, and any evaluation where the
stakes warrant human annotation.

\textbf{Mitigation:}  Swap response positions and check consistency, use
multiple judge models, validate against human annotations on a calibration set,
and always report which judge model was used.

\redflags
\begin{itemize}
    \item Treats LLM-as-judge as equivalent to human evaluation.
    \item Cannot name any specific bias.
    \item Proposes using the same model as both generator and judge.
    \item No awareness of position bias or how to mitigate it.
\end{itemize}

\interviewtest{Critical evaluation thinking, awareness of method limitations,
ability to reason about when a technique is and is not appropriate.}

\goodgreat
\begin{itemize}
    \item \badgeLfive{L5}: Names at least two biases, gives clear guidance on
          when to use vs.\ not.
    \item \badgeLsix{L6}: Describes mitigation strategies (position swapping,
          multi-judge ensemble, calibration), discusses the cost-quality tradeoff
          vs.\ human evaluation, and mentions specific tools (AlpacaEval,
          MT-Bench).
    \item \badgeLseven{L7}: Proposes a hierarchical evaluation pipeline
          (LLM-judge for development, human evaluation for deployment gates,
          production metrics for monitoring), discusses the meta-evaluation
          question (how do you evaluate the evaluator?), and connects to the
          broader philosophical issue of evaluation being the hardest unsolved
          problem in NLP.
\end{itemize}

\followups{How would you validate that your LLM judge agrees with humans?  What
is the inter-annotator agreement between GPT-4 and human evaluators?  How would
you handle a case where two LLM judges disagree?}
\end{interviewq}

\vspace{0.5em}

%% QUESTION 5 %%
\begin{interviewq}
\textbf{Q: Explain ROUGE. How is it different from BLEU?}

\badgeLfive{L5} \quad \badgetype{Conceptual} \quad \badgefreq{\freqhigh}

\stronganswer
ROUGE (Recall-Oriented Understudy for Gisting Evaluation) is the standard
automatic metric for summarization evaluation.  The key difference from BLEU
is the orientation:

\textbf{BLEU is precision-oriented:} ``Of the n-grams in the \emph{candidate},
how many appear in the reference?''  It penalizes candidates that add content
not in the reference.

\textbf{ROUGE is recall-oriented:} ``Of the n-grams in the \emph{reference},
how many appear in the candidate?''  It penalizes candidates that miss content
from the reference.

This difference is not arbitrary---it reflects the different failure modes of
each task.  In \textbf{translation}, the main risk is producing fluent but
unfaithful output (adding information not in the source), so
precision matters.  In \textbf{summarization}, the main risk is missing
important content, so recall matters.

\textbf{ROUGE variants:}
ROUGE-1 (unigram recall), ROUGE-2 (bigram recall), and ROUGE-L (longest common
subsequence, which captures in-order matching without requiring contiguity).

\textbf{Additional differences:}
(1)~BLEU uses a geometric mean across n-gram orders; ROUGE reports each variant
separately.
(2)~BLEU has a brevity penalty; ROUGE typically uses F-score variants.
(3)~BLEU is computed at corpus level; ROUGE is typically computed per-example
and averaged.

\textbf{Shared limitation:} Both are surface-form metrics that cannot capture
semantic equivalence or faithfulness.

\redflags
\begin{itemize}
    \item Confuses precision and recall orientation.
    \item Cannot explain why ROUGE is recall-oriented and BLEU is
          precision-oriented.
    \item Does not know the ROUGE variants (1, 2, L).
    \item Treats BLEU and ROUGE as interchangeable.
\end{itemize}

\interviewtest{Understanding of the design principles behind metrics and the
connection between metric design and task requirements.}

\goodgreat
\begin{itemize}
    \item \badgeLfive{L5}: Correctly states the precision vs.\ recall
          distinction, explains the three ROUGE variants, and gives the
          task-based rationale.
    \item \badgeLsix{L6}: Discusses ROUGE-LSum for multi-sentence evaluation,
          the faithfulness gap (ROUGE cannot detect hallucination), and
          BERTScore as a semantic improvement.
    \item \badgeLseven{L7}: Connects to the broader evaluation landscape
          (surface metrics $\to$ learned metrics $\to$ LLM-based evaluation),
          discusses ROUGE's bias toward extractive summaries, and explains
          why summarization evaluation remains an unsolved problem.
\end{itemize}

\followups{Why does ROUGE favor extractive summaries?  What is ROUGE-LSum?
How would you evaluate an abstractive summarizer where valid summaries may have
low ROUGE overlap with the reference?}
\end{interviewq}

\vspace{0.5em}

%% QUESTION 6 %%
\begin{interviewq}
\textbf{Q: Compare greedy decoding, beam search, and nucleus sampling. When
would you use each?}

\badgeLfive{L5} \quad \badgetype{Conceptual} \quad \badgefreq{\freqmed}

\stronganswer
These three decoding strategies represent fundamentally different tradeoffs
between quality, diversity, and computation:

\textbf{Greedy decoding} selects the highest-probability token at each step:
$w_t = \argmax P(w \mid w_{<t})$.  It is fast and deterministic but produces
repetitive, bland text because it always takes the locally optimal choice,
missing globally better sequences.  Use for: debugging, deterministic extraction
tasks, when diversity is harmful.

\textbf{Beam search} maintains $B$ candidate sequences and selects the
highest-scoring complete sequence.  It explores a wider space than greedy and is
excellent for tasks with a clear ``correct'' output, but it tends to produce
generic, safe text for open-ended tasks because high-likelihood sequences are
often repetitive or dull.  Use for: machine translation, speech recognition,
structured generation---tasks where there is a well-defined correct answer and
diversity is undesirable.

\textbf{Nucleus (top-$p$) sampling} draws from the smallest set of tokens whose
cumulative probability exceeds $p$ (e.g., 0.9).  Unlike top-$k$, the candidate
set size adapts to the model's confidence: few tokens when confident, many when
uncertain.  This produces diverse, natural-sounding text.  Use for: chatbots,
creative writing, any open-ended generation where diversity and naturalness
matter.

\textbf{The key design principle:}  Deterministic tasks (factual QA,
classification, extraction) benefit from low or no randomness.  Creative and
conversational tasks benefit from controlled randomness.  The decoding strategy
should match the task's tolerance for variation.

\redflags
\begin{itemize}
    \item Cannot explain the mechanism of any of the three strategies.
    \item Recommends beam search for open-ended chatbot generation.
    \item Does not explain why nucleus sampling is adaptive (unlike top-$k$).
    \item Says ``just use temperature 0.7'' without understanding the
          underlying strategy.
\end{itemize}

\interviewtest{Understanding of the generation process in autoregressive models,
the tradeoff between quality and diversity, and practical system design judgment.}

\goodgreat
\begin{itemize}
    \item \badgeLfive{L5}: Correctly describes all three mechanisms, gives the
          right use case for each, and explains the quality-diversity tradeoff.
    \item \badgeLsix{L6}: Discusses length normalization for beam search,
          temperature's interaction with nucleus sampling, repetition penalties,
          and how these strategies are combined in practice (e.g., top-$p$ +
          temperature + repetition penalty).
    \item \badgeLseven{L7}: Discusses speculative decoding for production
          serving, constrained decoding for structured output, the theoretical
          connection between beam search and MAP estimation, and why the
          highest-probability sequence is often not the best sequence for
          open-ended generation (the ``neural text degeneration'' finding).
\end{itemize}

\followups{What is the ``neural text degeneration'' problem?  How does
speculative decoding preserve the original distribution?  How would you choose
decoding parameters for a production chatbot?  What is constrained decoding?}
\end{interviewq}

\vspace{0.5em}

%% QUESTION 7 %%
\begin{interviewq}
\textbf{Q: How do you handle benchmark contamination? How would you design
an evaluation pipeline that is robust to it?}

\badgeLsix{L6} \quad \badgetype{Conceptual} \quad \badgefreq{\freqmed}

\stronganswer
Benchmark contamination occurs when benchmark questions or answers appear in
the model's pretraining data, inflating reported performance.  It is one of the
most significant threats to trustworthy LLM evaluation.

\textbf{Why it happens:}  LLMs are trained on massive web crawls (CommonCrawl,
etc.) that index pages containing benchmark datasets, their solutions, and
discussions about them.  Deduplication against known benchmarks helps but cannot
catch paraphrases, reformulations, or discussions that convey the answer
indirectly.

\textbf{Detection methods:}
\begin{enumerate}
    \item \textbf{N-gram overlap analysis:} Check for verbatim or near-verbatim
          overlap between training data and benchmark questions.  Limited by
          the availability of the training data (many models do not disclose it).
    \item \textbf{Canary strings:} Include unique identifier strings in benchmark
          data and check if the model can reproduce them.
    \item \textbf{Performance discontinuity:} Compare performance on contaminated
          vs.\ uncontaminated subsets.  Unexpectedly high performance on specific
          questions may indicate memorization.
    \item \textbf{Membership inference:} Test whether the model assigns
          suspiciously high likelihood to exact benchmark phrasings.
\end{enumerate}

\textbf{Robust evaluation design:}
\begin{enumerate}
    \item \textbf{Private held-out sets:} Create evaluation sets that are never
          published online.  This is the gold standard (used by Anthropic, Google
          internally).
    \item \textbf{Dynamic benchmarks:} Generate new questions regularly
          (LiveBench, SEAL).  Old questions are retired as contamination risk
          increases.
    \item \textbf{Multi-benchmark aggregation:} Report performance across
          many diverse benchmarks---contamination on one is diluted.
    \item \textbf{Human evaluation on novel prompts:} Chatbot Arena's approach
          of using real user prompts that did not exist during training.
    \item \textbf{Functional testing:} Instead of fixed question sets, test
          capabilities with procedurally generated problems (e.g., random
          arithmetic, novel code problems, shuffled answer options).
\end{enumerate}

\redflags
\begin{itemize}
    \item Not aware that contamination is a real and significant problem.
    \item Proposes only n-gram deduplication (which is necessary but insufficient).
    \item Cannot suggest any design principle for robust evaluation.
    \item Trusts reported benchmark numbers without questioning contamination.
\end{itemize}

\interviewtest{Evaluation maturity, awareness of the systemic challenges in LLM
evaluation, and ability to design robust evaluation systems.}

\goodgreat
\begin{itemize}
    \item \badgeLfive{L5}: Explains what contamination is and why it matters,
          suggests at least one detection method.
    \item \badgeLsix{L6}: Describes multiple detection methods and design
          principles (private held-out, dynamic benchmarks), discusses the
          tradeoff between benchmark standardization and contamination risk.
    \item \badgeLseven{L7}: Proposes a complete evaluation philosophy---static
          benchmarks for tracking, dynamic benchmarks for validation, private
          sets for deployment decisions, and production metrics for real-world
          assessment.  Connects to Goodhart's law and the broader challenge of
          measuring intelligence.
\end{itemize}

\followups{How does LiveBench address contamination?  What is the evidence that
contamination significantly inflates benchmark scores?  How would you build a
private evaluation set?  What is the role of functional testing?}
\end{interviewq}

\vspace{0.5em}

%% QUESTION 8 %%
\begin{interviewq}
\textbf{Q: Design an evaluation pipeline for a production chatbot that handles
customer support for an e-commerce company.}

\badgeLsix{L6} \quad \badgetype{System Design} \quad \badgefreq{\freqmed}

\stronganswer
I would design a three-tier evaluation pipeline: offline evaluation during
development, online evaluation at deployment, and continuous monitoring in
production.

\textbf{Tier 1: Offline evaluation (pre-deployment).}
\begin{itemize}
    \item \textbf{Regression test suite:} A curated set of 500+ representative
          customer queries with gold-standard responses, covering common intents
          (order status, returns, product questions) and edge cases (angry
          customers, ambiguous requests, out-of-scope questions).
    \item \textbf{Automated metrics:} BERTScore against gold responses for
          content quality, intent classification accuracy (does the chatbot
          identify the correct intent?), faithfulness checking (does the response
          contradict product/policy information?).
    \item \textbf{Safety testing:} Red-team prompts to test for prompt injection,
          information leakage (PII, internal policies), and harmful outputs.
    \item \textbf{LLM-as-judge:} GPT-4 evaluation on helpfulness, tone, and
          completeness for rapid iteration during development.
\end{itemize}

\textbf{Tier 2: Pre-launch human evaluation.}
\begin{itemize}
    \item \textbf{Expert evaluation:} Domain experts (customer support agents)
          evaluate 200+ conversations on accuracy, helpfulness, tone, and
          policy compliance.
    \item \textbf{Pairwise comparison:} Compare new model vs.\ current production
          model on the same queries.  Statistical test (Bradley-Terry or binomial)
          to confirm significance.
    \item \textbf{A/B test design:} Define primary metric (task completion rate),
          guardrail metrics (escalation rate, customer satisfaction score), and
          sample size calculation.
\end{itemize}

\textbf{Tier 3: Production monitoring.}
\begin{itemize}
    \item \textbf{Business metrics:} Task completion rate (did the customer's
          issue get resolved without human escalation?), average handling time,
          customer satisfaction (CSAT) score from post-chat survey.
    \item \textbf{Safety monitors:} Toxicity classifier on all outputs, PII
          detection in inputs and outputs, prompt injection detection.
    \item \textbf{Quality sampling:} Daily random sample of 50 conversations
          reviewed by QA team.  Track quality trends over time.
    \item \textbf{Drift detection:} Monitor query distribution changes, response
          latency, and failure rate.  Alert on anomalies.
    \item \textbf{User feedback:} Thumbs up/down on responses, escalation to
          human agent as implicit negative signal.
\end{itemize}

\textbf{Key design principles:}  Automated metrics catch regressions fast.
Human evaluation validates quality at deployment gates.  Production metrics
measure real-world impact.  No single tier is sufficient---all three are
necessary.

\redflags
\begin{itemize}
    \item Proposes only automated metrics with no human evaluation.
    \item No safety or guardrail testing.
    \item No production monitoring plan.
    \item Suggests a single metric (e.g., ``BLEU score'') as sufficient.
    \item Does not consider the business objective (task completion, not BLEU).
\end{itemize}

\interviewtest{End-to-end system design thinking, understanding of the
evaluation lifecycle (development $\to$ deployment $\to$ production), and
ability to connect technical metrics to business outcomes.}

\goodgreat
\begin{itemize}
    \item \badgeLfive{L5}: Proposes at least two tiers (automated + human),
          mentions safety testing and production monitoring.
    \item \badgeLsix{L6}: Designs all three tiers with specific metrics,
          includes A/B testing methodology, and connects to business KPIs
          (CSAT, escalation rate).
    \item \badgeLseven{L7}: Adds a continuous improvement loop (production
          failures feed back into the regression suite), discusses
          statistical rigor for A/B testing (power analysis, multiple
          comparisons correction), proposes a ``quality gate'' framework where
          deployment requires passing all tiers, and discusses the
          organizational challenges of evaluation (who owns the evaluation
          pipeline, how to handle disagreements between metrics and human
          judgment).
\end{itemize}

\followups{How would you handle a case where automated metrics pass but human
evaluators flag quality issues?  How do you set the sample size for A/B testing?
How would you detect if the chatbot is giving correct but unhelpful responses?
How do you evaluate multi-turn conversations where early errors compound?}
\end{interviewq}

\vspace{0.5em}

%% QUESTION 9 %%
\begin{interviewq}
\textbf{Q: You are building a code generation system. BLEU is commonly reported
for code generation---is it a good metric? What alternatives would you propose?}

\badgeLsix{L6} \quad \badgetype{Trade-off} \quad \badgefreq{\freqmed}

\stronganswer
BLEU is a \textbf{poor metric for code generation} despite being commonly
reported.  The fundamental problem: in code, surface form (tokens, formatting,
variable names) and \emph{functional correctness} are only loosely correlated.
Two programs can be functionally identical but share zero n-grams (different
variable names, different algorithms, different control flow).  Conversely, code
that is textually similar to the reference but has a single-character bug
(off-by-one, wrong operator) will have high BLEU but be completely wrong.

\textbf{Better alternatives:}

\begin{enumerate}
    \item \textbf{pass@k (HumanEval):} Generate $k$ candidate solutions, check
          how many pass all unit tests.  This directly measures functional
          correctness.  $\text{pass@}1$ is the most commonly reported.  The
          unbiased estimator (Chen et al., 2021) generates $n \geq k$ samples
          and computes: $\text{pass@}k = 1 - \binom{n-c}{k}/\binom{n}{k}$
          where $c$ is the number of correct samples.

    \item \textbf{Execution-based testing:} Run the generated code against a
          test suite and measure the fraction of tests that pass.  More general
          than pass@k (handles partial correctness).

    \item \textbf{CodeBERTScore:} Semantic similarity using CodeBERT embeddings.
          Better than BLEU for capturing semantic equivalence of code, but still
          does not guarantee functional correctness.

    \item \textbf{Edit similarity:} Tree-edit distance on the AST (abstract
          syntax tree).  Captures structural similarity while ignoring surface
          formatting.  Useful for measuring ``how close'' a wrong solution is.
\end{enumerate}

\textbf{The key insight:}  For code, the only ground truth is whether it
\emph{works}.  Execution-based metrics are strictly necessary.  Surface-form
metrics (BLEU, CodeBLEU, ROUGE) are useful auxiliary signals for measuring
diversity and similarity but should never be the primary metric.

\redflags
\begin{itemize}
    \item Accepts BLEU as appropriate for code generation without questioning it.
    \item Cannot articulate why surface-form metrics fail for code.
    \item Does not mention execution-based evaluation.
    \item Proposes ROUGE as a better alternative (it has the same fundamental
          problem).
\end{itemize}

\interviewtest{Critical evaluation thinking---can the candidate reason about
\emph{why} a metric is or is not appropriate for a given task, rather than
reflexively applying standard metrics?}

\goodgreat
\begin{itemize}
    \item \badgeLfive{L5}: Identifies that BLEU fails for code and proposes
          execution-based metrics.
    \item \badgeLsix{L6}: Explains pass@k with the unbiased estimator, discusses
          the limitations of execution-based metrics (test coverage,
          specification quality), and proposes a multi-metric approach.
    \item \badgeLseven{L7}: Discusses the broader challenge of evaluating
          generated code in production (security, efficiency, maintainability---
          not just correctness), formal verification approaches, and the
          connection to the program synthesis literature.
\end{itemize}

\followups{How do you handle the case where the test suite is incomplete?  What
is CodeBLEU and how does it differ from BLEU?  How would you evaluate code
generation for languages without good test infrastructure?}
\end{interviewq}

\vspace{0.5em}

%% QUESTION 10 %%
\begin{interviewq}
\textbf{Q: A model achieves state-of-the-art on MMLU but users report it gives
poor answers in practice. What could explain this gap?}

\badgeLsix{L6} \quad \badgetype{Trade-off} \quad \badgefreq{\freqmed}

\stronganswer
Several factors can cause a disconnect between benchmark performance and
real-world utility:

\textbf{1.\ Benchmark contamination:}  The model may have seen MMLU questions
(or very similar ones) during training, inflating its score without genuine
knowledge improvement.  This is the first thing I would investigate.

\textbf{2.\ Format mismatch:}  MMLU is multiple-choice (select from A/B/C/D).
Real users ask open-ended questions.  A model can be excellent at selecting
the correct option from a given set but poor at generating a correct, coherent
answer from scratch.  Multiple-choice benchmarks test recognition, not
generation.

\textbf{3.\ Distribution mismatch:}  MMLU covers 57 academic subjects.  Users
may be asking about topics not well represented (personal advice, domain-specific
jargon, current events).  High MMLU does not imply broad real-world knowledge.

\textbf{4.\ Missing evaluation dimensions:}  MMLU tests knowledge accuracy but
not helpfulness, instruction following, conversational ability, appropriate
hedging, tone, or safety.  A model can know the right answer but communicate it
poorly, refuse to answer unnecessarily, or give correct but unhelpful responses.

\textbf{5.\ Calibration and hedging:}  In multiple-choice, the model picks the
best option.  In open-ended generation, users want the model to express
uncertainty appropriately.  A model might be overconfident or underconfident in
ways that MMLU cannot measure.

\textbf{6.\ Alignment quality:}  MMLU measures pretrained knowledge.
Post-training alignment (SFT, RLHF) determines how that knowledge is surfaced
to users.  A model with excellent pretraining but poor alignment will score well
on MMLU but frustrate users.

\textbf{Lesson:}  No single benchmark captures real-world utility.  Production
evaluation must include diverse automatic benchmarks, human evaluation,
and real user feedback.

\redflags
\begin{itemize}
    \item Cannot hypothesize any explanation beyond ``the benchmark must be wrong.''
    \item Does not mention contamination.
    \item Treats high MMLU as proof of overall model quality.
    \item No awareness of the benchmark-to-deployment gap.
\end{itemize}

\interviewtest{Evaluation maturity, critical reasoning about metrics, and
understanding of the gap between benchmarks and real-world performance.}

\goodgreat
\begin{itemize}
    \item \badgeLfive{L5}: Identifies at least two plausible explanations
          (contamination, format mismatch).
    \item \badgeLsix{L6}: Provides a systematic analysis covering contamination,
          distribution mismatch, missing evaluation dimensions, and alignment
          quality.  Proposes concrete investigation steps.
    \item \badgeLseven{L7}: Connects to Goodhart's law, discusses how this
          phenomenon has played out historically in NLP (models gaming benchmarks
          from SQuAD to SuperGLUE to MMLU), proposes an evaluation philosophy
          that treats benchmarks as necessary but insufficient signals, and
          discusses the organizational implications (do not set MMLU as a team
          OKR).
\end{itemize}

\followups{How would you detect benchmark contamination in a model you did not
train?  What benchmarks would you add to MMLU to better predict user
satisfaction?  How would you design an evaluation framework that captures
real-world utility?}
\end{interviewq}

\bigskip

% ============================================================================
% CHAPTER SUMMARY
% ============================================================================
\section*{Chapter Summary}

\begin{enumerate}
    \item \textbf{Perplexity} ($\text{PP} = e^{H}$) is the standard language
          modeling metric, representing the weighted average branching factor.
          It measures prediction quality but not generation quality, faithfulness,
          or safety.  Use bits-per-byte for cross-tokenizer comparison.

    \item \textbf{BLEU} measures modified n-gram precision with a brevity penalty
          for machine translation.  It is precision-oriented, corpus-level, and
          has poor human correlation at the sentence level.  Neural metrics like
          COMET are superior but BLEU remains the reporting standard.

    \item \textbf{ROUGE} measures n-gram recall for summarization.  It is
          complementary to BLEU (recall vs.\ precision) but shares the same
          fundamental limitation: surface-form overlap does not capture semantic
          equivalence or faithfulness.

    \item \textbf{BERTScore} and \textbf{COMET} use learned embeddings to
          capture semantic similarity, achieving higher human correlation than
          surface-form metrics.

    \item \textbf{Classification metrics} (precision, recall, F1) must be
          chosen to match the task's error cost structure.  For NER, entity-level
          strict matching is the standard; token-level metrics overestimate
          performance.

    \item \textbf{Human evaluation} is the gold standard but is expensive, slow,
          and hard to reproduce.  Pairwise comparison is more reliable than
          Likert scales.  Chatbot Arena provides crowd-sourced evaluation at
          scale.

    \item \textbf{LLM-as-judge} scales better than human evaluation but has
          systematic biases (verbosity, position, self-preference).  It is
          a useful development tool but should not be the sole basis for
          deployment decisions.

    \item \textbf{Benchmarks} (MMLU, HumanEval, GSM8K) provide standardized
          comparison but suffer from contamination, saturation, and narrow
          coverage.  No single benchmark captures real-world utility.

    \item \textbf{Decoding strategies} trade off quality, diversity, and latency.
          Greedy and beam search for deterministic tasks; nucleus sampling for
          open-ended generation; constrained decoding for structured output;
          speculative decoding for production throughput.
\end{enumerate}

\begin{keyinsight}[Chapter 10 Summary: What Interviewers Are Really Testing]
When interviewers ask about evaluation, they are testing three things:
(1)~\emph{Do you know the math?}---Can you derive BLEU, define perplexity,
and explain why ROUGE is recall-oriented?
(2)~\emph{Do you know the limitations?}---Can you explain why BLEU fails for
paraphrases, why ROUGE favors extractive summaries, and why benchmark scores
do not guarantee real-world quality?
(3)~\emph{Can you design an evaluation pipeline?}---Can you propose a
multi-layered evaluation strategy that combines automatic metrics for speed,
human evaluation for quality validation, and production metrics for real-world
impact?  Candidates who demonstrate all three signal that they treat evaluation
as a first-class engineering problem, not an afterthought.
\end{keyinsight}
